% Section 2: setting, notation, targets, assumptions
\section{Setting, notation and assumptions}\label{sec:setting}

Source: PHASE1 \S\S1--3 (Appendix~\ref{app:phase1}) and PHASE2a \S0
(Appendix~\ref{app:phase2a}). Notation follows
\texttt{docs/extension/notation.tex}; Table~\ref{tab:notation} reproduces it.

\subsection{Data, latent rates and the accounting identity}

There are units $i=1,\dots,p$ (in Phase~3, $n$), nested in districts
$d=1,\dots,D$ (written $D^\star$ for the number of districts in PHASE2b and
PHASE3). Unit $i$ has size $N_i$, district $D_i$, row shares
$X_i\in\Delta^{R-1}$ and column shares $T_i\in\Delta^{C-1}$. The latent rate
$\beta_{irc}=N_{irc}/N_{ir}$ is the share of the $N_{ir}$ row-$r$ members of
unit $i$ who fall in column $c$, and $\beta_i\in\Bset:=(\Delta^{C-1})^R$. The
accounting identity is
\begin{equation}\label{eq:acc}
T_{ic}=\sum_{r=1}^RX_{ir}\beta_{irc},\qquad\text{i.e.}\quad T_i=\beta_i^\top X_i .
\end{equation}
The unit Duncan--Davis polytope $S(x,t)=\{b\in\Bset:b^\top x=t\}$ is the set of
rate matrices consistent with one unit's margins. Its support function is
$h_S(u;x,t)=\max_{b\in S(x,t)}\langle u,b\rangle$.

\subsection{Three targets}

PHASE1 \S1.7 and PHASE3 \S0.2 distinguish three targets. Each result says which
one it is about.
\begin{itemize}[nosep]
\item $B_{rc}=\sum_iN_iX_{ir}\beta_{irc}/\sum_iN_iX_{ir}$: the \emph{actual}
(realised, finite-population) rate in this electorate. This is JKST's target.
\item $\tilde B_{rc}=\sum_iw_{ir}m_{rc}(V_i)$, with $w_{ir}=N_iX_{ir}/\sum_jN_jX_{jr}$,
$V_i=(X_i,N_i,D_i)$ and $m=\E[\beta\mid V]$: the \emph{model-implied} rate at
the observed design. It is the conditional mean of $B$ given the design.
\item $B^*_{rc}=\E[NX_r\beta_{rc}]/\E[NX_r]$: the \emph{superpopulation} rate
used in the identification results.
\end{itemize}
$B=\tilde B+E$, where $E_{rc}=\sum_iw_{ir}e_{irc}$ and $e_i=\beta_i-m(V_i)$ is
the unit residual. Only $\tilde B$ and $B^*$ are affine in the unknown
parameter. All inference is conditional on the design and covers $B$ or
$\tilde B$; inference for $B^*$ is not pursued. Any linear target is written
$\varphi=\sum_{rc}a_{rc}B_{rc}$ (loyalty, defection, demobilisation, a gap
between groups).

\paragraph{In plain terms.} $B$ is ``what these voters actually did''.
$\tilde B$ is ``what the systematic pattern implies these units did''; it
ignores the local luck of each unit. The confidence interval for $B$ is wider
because it must allow for that luck.

\subsection{Linear contextual effects in homogeneous form}

Because $\mathbf 1^\top X=1$, an affine function of $X$ can be written as a
homogeneous linear one. The pooled LCE model (M0) says that for each column $c$
there is $H_c\in\mathbb R^{R\times R}$ with
\begin{equation}\label{eq:lce}
\E[\beta_{irc}\mid X_i,N_i]=e_r^\top H_cX_i\qquad\text{on }\{X_i\in A\},
\end{equation}
where $A\subseteq\Delta^{R-1}$ is the region on which the model is assumed
(by default the population hull $\conv\supp P_X$). Write
$H_c=Q_c+K_c$ with $Q_c=\tfrac12(H_c+H_c^\top)$ symmetric and
$K_c=\tfrac12(H_c-H_c^\top)$ skew. Then
\[
\tau_c(x):=\E[T_c\mid X=x]=x^\top H_cx=x^\top Q_cx,
\]
because $x^\top K_cx=0$. The regression identifies $Q$; it never sees $K$.
The unknown $K=(K_1,\dots,K_C)$ lives in
$\Kspace=\{K_c^\top=-K_c,\ \sum_cK_c=0\}$, of dimension $(C-1)R(R-1)/2$: one
number in $2\times2$, two in $2\times3$, six in $3\times3$.
$M_H(x)=[e_r^\top H_cx]_{rc}$ is the implied mean transition matrix.

\paragraph{What $H$ and $K$ mean (PHASE1 \S7.2).} $H_{c,rs}$ is the rate at which
row-$r$ voters go to column $c$ in a hypothetical unit made up entirely of
row-$s$ voters. $Q_{c,rs}$ is the average of ``$r$-voters in pure-$s$ units''
and ``$s$-voters in pure-$r$ units''; the data identify it. $K_{c,rs}$ is half
the gap between the two; the data never identify it. It is the R$\times$C version
of JKST's question ``is turnout low because Black residents vote less, or
because white residents in Black areas vote less?''

\subsection{Three identified sets and three hulls}

PHASE1 \S1.7 names three nested sets, $\ThI\subseteq\ThX\subseteq\ThO$
(PHASE2a Prop.~1):
\begin{itemize}[nosep]
\item $\ThI$: the sharp set, using the full law of $(N,X,T)$ and conditioning
on $(X,N)$ (PHASE2a Thm~1). Not a finite LP.
\item $\ThX$: the same, conditioning on $X$ only (JKST Prop.~1).
\item $\ThO$: the outer set, using only the regression $\E[T\mid X]$ and
probability validity (PHASE2a Prop.~7). A finite LP; sharp \emph{for the
regression function} only (PHASE2a Prop.~2).
\end{itemize}
The rule of the documents is that $\ThO$ is never called ``sharp'' without
``for the regression function''. There are also three hulls: the population hull
$A=\conv\supp P_X$ (the default in identification), the hull of the empirical
law (exact for finite-sample identification), and the sample hull used as an
estimate of $A$. Since the sample hull lies inside $A$, estimating it only
widens the outer bound (PHASE3 Prop.~P3.2(c)). Every ``sharp'' statement is
relative to continuous (real-valued) latent rates, not integer cell counts
(PHASE2a Remark~4.2).

\subsection{The assumptions, and what each implies}\label{sec:assumptions}

Each assumption is stated as in PHASE1 \S2, followed by a plain-language
reading: what it rules in or out, when it is credible, and what goes wrong if it
fails. Which rungs use which assumption is shown in Table~\ref{tab:assumption-use}.

\begin{assumption}[A1: closed population / accounting identity]
$T_i=\beta_i^\top X_i$ exactly, with $N_i$ the same people in both margins.
\end{assumption}
\begin{plain}
Both margins describe the same people: the same electorate in both elections,
or the same eligible population for group and outcome. It rules out turnover,
moving, deaths, new registrants and boundary changes. It is credible when both
margins come from one register and the row and column categories include
abstention and, where relevant, ``entered'' or ``left''. If it fails, the
Duncan--Davis sets and the regression identity both shift, and every bound can
exclude the truth. Bounded turnover is an open extension: it cannot be handled
by widening the final interval (PHASE1 A1, as corrected after review). Every
rung, including DD, needs A1.
\end{plain}

\begin{assumption}[A2: linear contextual effects on $A$]
Equation~\eqref{eq:lce}: each group's mean rate is linear in the unit's
composition on $A$.
\end{assumption}
\begin{plain}
A group's behaviour may depend on its surroundings, but only along a straight
line: for example, in-party loyalty rises steadily as the in-party share of the
poll rises. It rules out thresholds and tipping (a machine that mobilises only
above 50\%) and U-shapes (competitive polls mobilising both sides). It is the
\emph{only} identifying restriction of the LCE rungs. If it fails, the
quadratic regression is misspecified and the bound can miss the truth by an
amount that does not shrink with more data. Some violations are invisible:
nonlinear departures that cancel in the aggregate leave the data unchanged
while moving the target (PHASE2a Prop.~9). A cubic term or an empty bound can
reveal some failures; goodness of fit can never validate A2.
\end{plain}

\begin{assumption}[A3: mean independence from size]
$\E[\beta\mid X,N]=\E[\beta\mid X]$ (with $D$ added to both sides under the
district rungs).
\end{assumption}
\begin{plain}
Large and small units with the same mix behave alike. It matters because the
target weights units by size. If large urban polls differ in loyalty from small
rural ones with the same partisan mix, the size-weighted target is biased by
$\Cov(NX_r,\beta_{rc}-M(X))/\E[NX_r]$. District dummies help when size differs
by district (PHASE2b Prop.~D4). It is partly testable through $X\cdot\log N$
interactions.
\end{plain}

\begin{assumption}[A4: support and extrapolation]
$\supp(X)\subseteq A$, compact, and A2 holds on all of $A$.
\end{assumption}
\begin{plain}
The straight lines are trusted over the whole region $A$, whose corners (the
most extreme units) set the constraints. A larger $A$ gives tighter bounds but
assumes more: using the whole simplex assumes the lines extend to pure units
that may not exist (PHASE2a Prop.~5). Outlier units at the hull's corners can
make constraints bind spuriously. Check by trimming the hull. PHASE2a adds
the condition (S3), probability validity on all of $A$; it is automatic for
the population hull (PHASE2a Lemma~3(c)).
\end{plain}

\begin{assumption}[A5: rank / richness]
For symmetric $D$, $x^\top Dx=0$ for all $x\in\supp_AP_X$ implies $D=0$, where
$\supp_AP_X$ is the support of $P_X$ restricted to $A$.
\end{assumption}
\begin{plain}
The compositions vary enough to fit the quadratic regression. For two groups it
means at least three distinct compositions. For three groups it needs six
compositions not on a common conic: five are never enough (PHASE2a
Remark~1.1). It fails, or nearly fails, when shares barely vary (JKST's
``gender'' problem). Then $Q$ is weakly identified and the bound is wide and
unstable. Check the condition number of the design matrix.
\end{plain}

\begin{assumption}[A6: source of randomness and independence]
Given the design $\{(N_i,X_i,D_i)\}$ and the district effects, the unit rates
$\beta_i$ are random with mean $m(V_i)$, and the residuals $e_i$ are
independent across units (P3-S1a).
\end{assumption}
\begin{plain}
Every unit is observed, so the uncertainty is not from sampling units. It is
model-based: given its characteristics, each unit's behaviour has a random part.
District shocks are fixed features of this electorate (part of the mean under
M1/M2; a misspecification under M0). Independence is a substantive assumption.
Neighbourhood spillovers within a district violate it, and then the variances
and residual allowances are too small. Observing every unit does \emph{not}
remove dependence (PHASE3 \S0.4, as corrected in round~2). The pooling test
does not test it.
\end{plain}

\begin{assumption}[A7: regularity for inference]
A unique active constraint at the optimum, $\mu_{rs}>0$, and $\ThO$ not a
singleton.
\end{assumption}
\begin{plain}
This is JKST's tie-breaking condition. With it the bound's endpoint is a smooth
function of the regression coefficients and the delta method applies. In
R$\times$C, and especially with many districts, near-ties are the norm, so the
inference of Phase~3 is built not to need it.
\end{plain}

\begin{assumption}[A8: homogeneous contextual effects across districts (M1)]
$K$ (and the symmetric slope parts) are common to all districts; only
intercepts vary.
\end{assumption}
\begin{plain}
The way composition shifts behaviour is the same in every district, while each
district may be more or less loyal across the board (parallel lines). It is
what produces the intersection gain of M1. It fails if, say, a strong
incumbent changes the contextual gradient in one district, or if ceiling and
floor effects make an additive shift impossible near 0 or 1, which is most
likely in the lopsided districts that tighten the bound most. The bound can then
exclude the truth and still look narrow. The symmetric part is testable
(PHASE2b Lemma~D1(b)); the shared $K$ is not.
\end{plain}

\begin{assumption}[A9: within-district richness]
Under M2, A5 within each district. Under M1, each district's support contains
$R$ affinely independent compositions and the stacked design has full column
rank (PHASE2b Lemma~D1(d)).
\end{assumption}
\begin{plain}
Districts with few, similar units identify almost nothing. Under M1 the bar is
lower than under M2: a district's baseline needs only $R$ affinely
independent compositions. Weak identification per district, combined with a
max/min over districts, biases plug-in bounds inward.
\end{plain}

\begin{assumption}[A10: asymptotic regime (a choice)]
(R1) $D$ fixed, $p_d\to\infty$; (R2) $D\to\infty$, $p_d$ bounded; (R3) both
grow.
\end{assumption}
\begin{plain}
Which limit theory applies depends on the design. With ten polls in each of
ninety ridings (R2), district intercepts are incidental parameters and the
plug-in intersection is badly too narrow. PHASE3 \S0.5 says which result holds
in which regime (Table~\ref{tab:regimes} in Section~\ref{sec:inference}).
\end{plain}

\paragraph{Assumptions explicitly not made.} No distributional form for $\beta$
(unlike King's truncated normal), no independence between $\beta$ and $X$
(contextual effects are the point), and no point-identifying nonlinearity.

\begin{longtable}{@{}p{0.30\linewidth}p{0.62\linewidth}@{}}
\caption{Which rungs use which assumption (PLAIN\_LANGUAGE\_GUIDE \S2; PHASE2c \S0).}\label{tab:assumption-use}\\
\toprule Rung & Assumptions used \\ \midrule\endfirsthead
\toprule Rung & Assumptions used \\ \midrule\endhead
DD & A1 only. \\
N1 (sign rows on realised tables) & A1; continuous unit tables; no regression and no sampling inference. \\
M0 & A1--A5; inference adds A6, and A7 only for the normal limit of P3.4(i). \\
M1, M1($\delta$), M2, L4, L5, L6 & A1--A5 within districts (A5 replaced by A9), A8 for M1 and its relatives; inference adds A6 and the regime A10. \\
N2 & A1, A3 and the finite-support convention of PHASE2c \S4. \\
M-SIGN rows & Those of the rung they are added to. \\
\bottomrule
\end{longtable}

\begin{table}[ht]
\centering\small
\begin{tabular}{@{}ll@{}}
\toprule
Symbol & Meaning \\
\midrule
$i=1,\dots,p$ & unit (precinct) index \\
$d=1,\dots,D$;\ $D_i$;\ $\mathcal I_d$;\ $p_d$ & district; district of unit $i$; units in $d$; their number \\
$r,s=1,\dots,R$ & row categories (election-1 choice / group) \\
$c=1,\dots,C$ & column categories (election-2 choice) \\
$N_i,\ N_{ir},\ N_{irc}$ & unit size; row count; latent cell count \\
$X_i\in\Delta^{R-1}$;\ $T_i\in\Delta^{C-1}$ & observed row shares; observed column shares \\
$\beta_i=[\beta_{irc}]\in\mathfrak B=(\Delta^{C-1})^R$ & latent transition matrix of unit $i$ \\
$W_i=(N_i,X_i,T_i,Z_i,D_i)$;\ $V=(X,N,D)$ & observables of unit $i$; conditioning variables \\
$P$;\ $\mathbb Q$ & law of $W$ (observable); law of $(W,\beta)$ \\
$\mathcal G=\sigma(V)$ & conditioning $\sigma$-algebra of the contextual model \\
$B_{rc}$,\ $\tilde B_{rc}$,\ $B^*_{rc}$ & actual, model-implied, superpopulation rate \\
$\varphi=\sum_{rc}a_{rc}B_{rc}$ & linear target \\
$A\subseteq\Delta^{R-1}$,\ $\operatorname{vert}(A)$ & region where LCE is assumed; its vertices \\
$H_c$;\ $Q_c=\tfrac12(H_c+H_c^\top)$;\ $K_c=\tfrac12(H_c-H_c^\top)$ & LCE coefficients; identified part; unidentified part \\
$\mathcal K$;\ $k$ & space of $K$, $\dim=(C-1)R(R-1)/2$; its free coordinates \\
$\kappa_c=K_{c,12}$ & unidentified scalar per column when $R=2$ \\
$M_H(x)=[e_r^\top H_cx]_{rc}$;\ $\tau_c(x)=E[T_c\mid X=x]$ & implied mean transition matrix; regression function \\
$\mu_r=E[NX_r]$,\ $\mu_{rs}=E[NX_rX_s]$ & weighting moments \\
$S(x,t)$;\ $h_S(u;x,t)$ & unit Duncan--Davis polytope; its support function \\
$\Theta_I\subseteq\Theta_X\subseteq\Theta_O$ & sharp set; JKST Prop.~1 set; outer set \\
$\zeta_{c,d}\in\mathbb R^R$ & district intercept shifts (M1), identified \\
$Q^{(d)}_c,\ K^{(d)}_c,\ A_d$ & district-level identified part, unidentified part, support hull \\
$\delta_B,\ \delta_W,\ \Delta$ & between, within (W1) and pooled-within (W2) associations (PHASE2c) \\
$\mathrm{CI}_x,\ S_1$ & JKST confidence interval; residual (Popoviciu) standard-error term \\
$\Phi$ & standard normal CDF \\
\bottomrule
\end{tabular}
\caption{Notation (after \texttt{docs/extension/notation.tex}, extended with symbols used later).}\label{tab:notation}
\end{table}
